\documentclass[10pt,a4paper]{amsart}
\usepackage[margin=19mm]{geometry}
\usepackage{amsmath,amssymb,mathtools}
\usepackage[scr=rsfs,cal=euler]{mathalfa}
\usepackage{xfrac}
\usepackage[unicode,hidelinks,bookmarks=false]{hyperref}
\newcommand{\ie}{i.e.\ }
\newcommand{\cf}{cf.\ }
\newcommand{\resp}{respectively\ }
\newcommand{\Subsection}{\subsection}
\providecommand{\nobreakdash}{\nobreak\hbox{-}}
\setlength{\emergencystretch}{2em}
\multlinegap0pt
\usepackage{dsfont}
\usepackage{amssymb}
\usepackage{tikz}
\usepackage{enumerate}
\usepackage[shortlabels]{enumitem}
\usepackage{xcolor}
\usepackage[normalem]{ulem}
\newtheorem{prop}{Proposition}
\newcommand{\PP}{\mathds{P}}
\newcommand{\PPP}{\PP^{3}}
\title{The classification of ACM curves on a surface in $\PPP$}
\author{Abel Castorena, Montserrat Vite}
\date{Source: arXiv:2602.06141v1 (CC BY 4.0)}
\begin{document}
\maketitle

\noindent\emph{Source and licence.} The classification of ACM curves on a surface in $\PPP$. Version arXiv:2602.06141v1 (CC BY 4.0), distributed by arXiv under the Creative Commons Attribution 4.0 International licence, http://creativecommons.org/licenses/by/4.0/, as recorded in the arXiv licence metadata for this exact version and shown on the abstract page https://arxiv.org/abs/2602.06141v1. Full licence notice: \texttt{licenses/arxiv-2602.06141-CC-BY-4.0.txt}.\par\smallskip

\noindent\textit{Editorial scope.} Counted unit: the article's prop propnotdet (source0.tex lines 262--264). The article's complete original proof of it is delivered with it (source0.tex lines 266--291, one \texttt{proof} environment of 26 lines). No mathematics is written, repaired or re-derived: the statement and the proof are the article's own lines, in the article's own notation, and no retained line is substituted or re-flowed.  No further source-local context is needed: every cross-reference of every retained line resolves inside the two delivered spans.  Nothing was omitted from any retained span, so the package carries no omission record.  No retained line contains a citation, so the wrapper carries no bibliography and no bibliography entry.  No retained line contains a figure environment, an image-inclusion command or drawing code, so no image is delivered and the package declares no asset dependency.  Editorial formatting only, in addition: an AMS \texttt{amsart} preamble that restates the article's own declarations and macros used by the retained lines, namely the environments \texttt{prop} on separate counters, together with the standard packages those lines need.  The retained environments print in this document as Proposition 1 in order of appearance, whereas the article prints the same items with the numbering of its own full document.  The complete native source of the article as distributed in the arXiv e-print for this exact version is bundled unchanged as \texttt{sources/194080/source0.tex}.
\begin{prop} \label{propnotdet}
Let $X=V(F)\subseteq \PPP$ be a surface of degree $d$. If there exists an ACM curve $C\subseteq X$ that is not a complete intersection, then $X$ is a weak determinantal surface of type $(\hat{a}, \hat{b})$ for some weak admissible pair $(\hat{a}, \hat{b})$ of degree $d$ and length $t$.
\end{prop}

\begin{proof}
Let $C\subseteq X$ be an ACM curve that is not a complete intersection. Two cases may occur.

\begin{itemize}
\item[Case 1:] $F$ is a minimal generator of the ideal of $C$.\\
By the Hilbert--Burch Theorem, there exists a matrix $M$ of size $(t+1)\times t$ such that $F$ is the determinant of a $t\times t$ minor of $M$. By degree considerations, we have $t\leq d$. After a change of coordinates, we may assume that the Betti numbers of the minimal free resolution of the ideal of $C$ are ordered such that $a_{0}=d$, $a_{1}\leq\cdots\leq a_{t}$, and $b_{1}\leq\cdots\leq b_{t}$. Let $m=(m_{ij})$ be the $t\times t$ submatrix of $M$ such that $\det(m)=F$ and $\deg(m_{ij})=b_j-a_i$. Setting $\hat{a}=(a_1,\ldots,a_t)$ and $\hat{b}=(b_1,\ldots,b_t)$, it follows that $X$ is weak determinantal of type $(\hat{a},\hat{b})$.

\item[Case 2:] $F$ is not a minimal generator of the ideal of $C$.\\
Let $F_1,\ldots,F_t$ be a minimal generating set of the ideal of $C$. By the Hilbert--Burch Theorem, there exists a matrix $M=(m_{ij})$ of size $t\times(t-1)$ such that each $F_i$ is the determinant of a $(t-1)\times(t-1)$ minor of $M$. Since $F$ lies in the ideal of $C$, there exist homogeneous polynomials $m_{1t},\ldots,m_{tt}\in\mathds{C}[x,y,z,w]$ such that
\[
F=\sum_{i=1}^{t}(-1)^i m_{it}F_i.
\]
Let $S$ be the matrix obtained by completing $M$ with the column $(m_{it})$. Expanding $\det(S)$ along the last column yields $\det(S)=\pm F$. The degrees satisfy $\deg(m_{ij})=b_j-a_i$, where $a_1,\ldots,a_t$ and $b_1,\ldots,b_{t-1}$ are the Betti numbers of the minimal free resolution of the ideal of $C$, ordered increasing. Setting $b_t=d$, $\hat{a}=(a_1,\ldots,a_t)$, and
\[
\hat{b}=(b_1,\ldots,b_{j_0-1},b_t,b_{j_0},\ldots,b_{t-1})
\quad\text{if } b_{j_0-1}\leq b_t<b_{j_0},
\]
we conclude that $X$ is a weak determinantal surface of type $(\hat{a},\hat{b})$.

Finally, we must verify that $m_{it}=0$ whenever $a_i=d$. If not, there exists $F_i$ of degree $d$ such that $m_{it}\neq0$, and hence
\[
F_i=\sum_{j\neq i} m_{jt}'F_j + m_{it}^{-1}F.
\]
Thus, $\{F_1,\ldots,F_{i-1},F,F_{i+1},\ldots,F_t\}$ would be a minimal generating set of the ideal of $C$, contradicting the hypothesis.
\end{itemize}
\end{proof}

\end{document}
